\documentclass{beamer}
\usepackage[utf8]{inputenc}

\usetheme{Madrid}
\usecolortheme{default}
\usepackage{amsmath,amssymb,amsfonts,amsthm}
\usepackage{txfonts}
\usepackage{tkz-euclide}
\usepackage{listings}
\usepackage{adjustbox}
\usepackage{array}
\usepackage{tabularx}
\usepackage{gvv}
\usepackage{lmodern}
\usepackage{circuitikz}
\usepackage{tikz}
\usepackage{graphicx}

\setbeamertemplate{page number in head/foot}[totalframenumber]

\usepackage{tcolorbox}
\tcbuselibrary{minted,breakable,xparse,skins}



\definecolor{bg}{gray}{0.95}
\DeclareTCBListing{mintedbox}{O{}m!O{}}{%
	breakable=true,
	listing engine=minted,
	listing only,
	minted language=#2,
	minted style=default,
	minted options={%
		linenos,
		gobble=0,
		breaklines=true,
		breakafter=,,
		fontsize=\small,
		numbersep=8pt,
		#1},
	boxsep=0pt,
	left skip=0pt,
	right skip=0pt,
	left=25pt,
	right=0pt,
	top=3pt,
	bottom=3pt,
	arc=5pt,
	leftrule=0pt,
	rightrule=0pt,
	bottomrule=2pt,
	toprule=2pt,
	colback=bg,
	colframe=orange!70,
	enhanced,
	overlay={%
		\begin{tcbclipinterior}
			\fill[orange!20!white] (frame.south west) rectangle ([xshift=20pt]frame.north west);
	\end{tcbclipinterior}},
	#3,
}
\lstset{
	language=C,
	basicstyle=\ttfamily\small,
	keywordstyle=\color{blue},
	stringstyle=\color{orange},
	commentstyle=\color{green!60!black},
	numbers=left,
	numberstyle=\tiny\color{gray},
	breaklines=true,
	showstringspaces=false,
}
%------------------------------------------------------------
%This block of code defines the information to appear in the
%Title page
\title %optional
{5.2.49}
%\subtitle{A short story}

\author % (optional)
{RAVULA SHASHANK REDDY - EE25BTECH11047}

 \begin{document}
	
	
	\frame{\titlepage}
	\begin{frame}{Question}
     Consider the basis 
\[
\{\vec{u}_1, \vec{u}_2, \vec{u}_3\} \text{ of } \mathbb{R}^3, \text{ where }
\vec{u}_1 = \myvec{1\\0\\0}, \ 
\vec{u}_2 = \myvec{1\\1\\0}, \
\vec{u}_3 = \myvec{1\\1\\1}.
\]
Let $\{\vec{f}_1, \vec{f}_2, \vec{f}_3\}$ be the dual basis and 
$f(a,b,c) = a+b+c$. If 
\[
\vec{f} = \alpha_1 \vec{f}_1 + \alpha_2 \vec{f}_2 + \alpha_3 \vec{f}_3,
\] 
find $\myvec{\alpha_1& \alpha_2 &\alpha_3}$.
\end{frame}
\begin{frame}{Solution}
    The dual basis satisfies :
\begin{align}
\vec{f}_i(\vec{u}_j) &= \delta_{ij} = 
\begin{cases} 
1, & i=j \\ 
0, & i\neq j 
\end{cases}. \label{eq:delta}
\end{align}

\begin{align}
\vec{U} &= \myvec{\vec{u}_1 \ \vec{u}_2 \ \vec{u}_3} 
= \myvec{1 & 1 & 1\\ 0 & 1 & 1\\ 0 & 0 & 1}, \\
\vec{F} &= \myvec{\vec{f}_1 \\ \vec{f}_2 \\ \vec{f}_3}.
\end{align}

Then the dual basis condition becomes
\begin{align}
\vec{F} \, \vec{U} &= \vec{I_3}. \label{eq:F_U}\\
\vec{F}^{-1} &= \vec{U} 
\end{align}
\end{frame}
\begin{frame}{Solution}

\begin{align}
\vec{f} &= \alpha_1 \vec{f}_1 + \alpha_2 \vec{f}_2 + \alpha_3 \vec{f}_3 
\implies \vec{f} = \vec{a}^T \vec{F}, \\
\text{where}\quad\vec{a}&=\myvec{\alpha_1\\ \alpha_2\\\alpha_3} \quad \vec{f}=\myvec{1&1&1}  \\
\vec{a}^T \vec{F} &= \vec{f} \implies \vec{a}^T = \vec{f} \vec{F}^{-1} = \vec{f} \vec{U}, \\
\vec{a}^T &= \myvec{1 & 1 & 1} \myvec{1 & 1 & 1\\ 0 & 1 & 1\\ 0 & 0 & 1} 
= \myvec{1 & 2 & 3}.
\end{align}

\begin{align}
\myvec{\alpha_1\\ \alpha_2\\\alpha_3} = \myvec{1\\2\\3}
\end{align}
\end{frame}
\begin{frame}[fragile]
\frametitle{C Code}
\begin{lstlisting}
    #include <stdio.h>

int main(void) {
    // Basis matrix U: columns are u1,u2,u3 but stored row-major U[row][col]
    double U[3][3] = {
        {1.0, 1.0, 1.0},
        {0.0, 1.0, 1.0},
        {0.0, 0.0, 1.0}
    };

    // f as row vector (f = (1,1,1))
    double f[3] = {1.0, 1.0, 1.0};

    // alpha will store the row vector alpha = f * U
    double alpha[3] = {0.0, 0.0, 0.0};

    // Compute alpha_j = sum_k f_k * U[k][j]
    \end{lstlisting}
    \end{frame}
\begin{frame}[fragile]
\frametitle{C Code}
\begin{lstlisting}
    for (int j = 0; j < 3; ++j) {
        double sum = 0.0;
        for (int k = 0; k < 3; ++k) {
            sum += f[k] * U[k][j];
        }
        alpha[j] = sum;
    }

    // Print result
    printf("alpha = (");
    for (int j = 0; j < 3; ++j) {
        printf("%.0f", alpha[j]);            // integer values expected, print without decimals
        if (j < 2) printf(", ");
    }
    printf(")\n");

    return 0;
}
\end{lstlisting}
\end{frame}
\begin{frame}[fragile]
\frametitle{Python Direct}
\begin{lstlisting}
    import numpy as np

# Given basis matrix U
U = np.array([
    [1, 1, 1],
    [0, 1, 1],
    [0, 0, 1]
], dtype=float)

# Given linear functional f(a,b,c) = a + b + c
f = np.array([1, 1, 1], dtype=float)

# Compute alpha^T = f * U
alpha = f @ U   # matrix multiplication

print("alpha =", alpha)
\end{lstlisting}
\end{frame}
\begin{frame}[fragile]
\frametitle{Python Shared}
\begin{lstlisting}
    import ctypes
import numpy as np

# Load the shared library
lib = ctypes.CDLL("./libalpha.so")

# Define argument and return types
lib.compute_alpha.argtypes = [
    (ctypes.c_double * 3 * 3),  # 3x3 matrix
    (ctypes.c_double * 3),      # f vector
    (ctypes.c_double * 3)       # output alpha
]

# Prepare data
U = np.array([
    [1.0, 1.0, 1.0],
    [0.0, 1.0, 1.0],
    [0.0, 0.0, 1.0]
    \end{lstlisting}
\end{frame}
\begin{frame}[fragile]
\frametitle{Python Shared}
\begin{lstlisting}
], dtype=np.double)

f = np.array([1.0, 1.0, 1.0], dtype=np.double)
alpha = np.zeros(3, dtype=np.double)

# Convert numpy arrays to ctypes
U_ctypes = (ctypes.c_double * 3 * 3)(*map(lambda row: (ctypes.c_double * 3)(*row), U))
f_ctypes = (ctypes.c_double * 3)(*f)
alpha_ctypes = (ctypes.c_double * 3)(*alpha)

# Call C function
lib.compute_alpha(U_ctypes, f_ctypes, alpha_ctypes)

# Convert result back to numpy array
alpha_result = np.array(list(alpha_ctypes))

print("alpha =", alpha_result)
\end{lstlisting}
\end{frame}
\end{document}
